\chapter*{Remerciements}
\phantomsection
\addcontentsline{toc}{chapter}{Remerciements}

Je tiens à exprimer ma profonde gratitude à l'École Marocaine des Sciences de l'Ingénieur (EMSI Rabat) pour la qualité de la formation d'ingénierie en informatique et réseaux, ainsi que pour l'accompagnement méthodologique qui a structuré la démarche scientifique et technique de ce mémoire.

Je remercie chaleureusement mon encadrant pédagogique à l'EMSI, \textbf{M. Khalid ZINE DINE}, pour la disponibilité, les orientations pertinentes et l'exigence constructive portée sur la cohérence du rapport, la validation des choix d'architecture et la qualité de la rédaction.

Mes remerciements vont également à l'agence d'accueil \textbf{Agence de Développement du Digital (\#ADD)} et à mon encadrant de stage, \textbf{M. Wail DARDAR}, pour la confiance accordée, la définition progressive des objectifs, et la transmission d'une vision réaliste des contraintes industrielles (délais, priorisation, qualité de livraison, sécurité).

Je souhaite saluer l'équipe de stage pour l'esprit collaboratif, l'entraide technique et l'enrichissement professionnel qu'elle a apportés tout au long de la période : \textbf{Yasmine EL HABBOUR}, \textbf{MRANI ZENTAR Inaam}, \textbf{CHIKI Hind}, et \textbf{Amine EL MEZOUAGHI}. Les échanges réguliers, les revues de solutions et le partage d'expérience ont contribué directement à la maturation des choix techniques et à la validation empirique de plusieurs hypothèses de conception.

Je remercie enfin l'ensemble du corps enseignant et mes camarades de promotion pour les discussions enrichissantes, ainsi que ma famille pour le soutien indispensable durant les phases intenses de rédaction et de finalisation.

\vspace{1cm}
\hfill Asaad AZAROUAL

\cleardoublepage